\section{Related work}

\textbf{Parameter-Efficient Fine-Tuning (PEFT)} methods are designed to reduce the high expense of fine-tuning large-scale models. They achieve this by training a relatively small subset of parameters, compared to the total number of parameters, for adapting to downstream tasks. Existing PEFT strategies can be divided into two categories: \textit{Prompt-based} methods add extra soft tokens (prompts) to the initial input and focus solely on fine-tuning these trainable vectors~\cite{26lester2021prompttuning, Razdaibiedina2023progressiveprompts, wang2021learningtoprompt}. \textit{Adapter-based methods} introduce additional trainable modules into the original frozen backbone~\cite{13houlsby2019peft, 22rebuffi2017residualadapters, 23pfeiffer2020adapterfusion, 24ruckle2020adapterdrop}. LoRA~\cite{1hu2021lora} expands upon adapter-based fine-tuning by adding a small number of trainable low-rank matrices alongside layers of frozen weights, which introduces a negligible inference overhead. Variants of LoRA include works like~\cite{meng2024periodiclora}, which employs SVD decomposition to prune less significant singular values for more efficient updates. Another variation, DoRA ~\cite{dora}, decomposes pre-trained weights into magnitude and direction components while applying LoRA the latter. QLoRA~\cite{11dettmers2023qlora} optimizes LoRA’s design one step further, using 4-bit NF4 weights, double quantization to reduce the memory footprint, and paged optimizers to alleviate memory spikes. In our experiments, we focus on the original implementation of LoRA with 4-bit quantization.

\textbf{Efficient serving of LoRA models.} The main challenges for serving multiple fine-tuned models efficiently are:

\begin{enumerate}
\item \textbf{Scalability:} As the demand for model inference grows, the system must scale efficiently to handle the increased load. This involves not just scaling up the computational resources but also managing the load distribution among models to maintain performance.

\item \textbf{Cost:} The computational resources required to serve multiple fine-tuned models can lead to significant costs. Efficiently managing these costs while maintaining high performance and availability is a major challenge.

\end{enumerate}

Techniques like Segmented Gather Matrix-Vector Multiplication (SGMV)~\cite{38chen2023punica} aim to address these challenges by optimizing the way computations are performed and resources are used. Open source tools like DeepSpeed\footnote{https://github.com/microsoft/DeepSpeed}, FasterTransformer\footnote{https://github.com/NVIDIA/FasterTransformer}, and vLLM~\cite{41kwon2023efficient} also aim to enable cost-effective and scalable serving of fine-tuned models. In this paper, we use LoRAX\footnote{https://github.com/predibase/lorax}, which is specifically designed for the efficient serving of LLMs fine-tuned with LoRA. LoRAX supports dynamic adapter loading so adapters can be downloaded asynchronously during inference, multiple model families like Llama~\cite{4llama} and Mistral~\cite{5mistral}, and bitsandbytes\footnote{https://github.com/TimDettmers/bitsandbytes}-quantized models.
